% TOP_PROBLEM: {"id": 3120, "title": "Matching cut and girth", "category_id": 3, "category": {"id": 3, "name": "graph_theory", "display_name": "Graph Theory", "description": "Problems involving graphs, networks, and their properties.", "slug": "graph-theory", "order_index": 3, "created_at": "2026-07-31T15:26:25.671Z"}, "status": "open"}
% FIRST_POSTED: null
% ENTRY_KIND: statement_only
% Imported from: batch02.tex; UnsolvedMath ID 3120
% Compile twice: lualatex 3120.tex
\documentclass[11pt]{article}
\usepackage{fontspec}
\setmainfont{Latin Modern Roman}
\usepackage{amsmath,amssymb,mathtools,unicode-math}
\setmathfont{Latin Modern Math}
\usepackage{xurl,hyperref}
\hypersetup{hidelinks,pdftitle={UnsolvedMath Problem 3120}}
\newcommand{\sref}[2]{\hyperlink{source:#1:#2}{[S#2]}}
\newcommand{\cataloguescope}[1]{\paragraph{Mathematical question.}#1}
\begin{document}
% UnsolvedMath ID 3120; source catalogue code OPG-37364
\setcounter{section}{3119}
\section{Matching cut and girth}\label{3120}
{\small\sffamily UnsolvedMath ID 3120 · Graph Theory}
\subsection{Definitions and mathematical statement}

\paragraph{Shared notation.}$\mathbb N=\{1,2,\ldots\}$, and $\mathbb Z,\mathbb Q,\mathbb R,\mathbb C$ denote the integers, rationals, reals and complex numbers. Any other conventions are specified locally.


For a finite simple graph $G=(V,E)$, a matching is an edge set whose edges have pairwise disjoint endpoints. A matching cut is a set $\delta_G(S)=\{uv\in E:u\in S,v\in V\setminus S\}$ that is a matching, where $\varnothing\ne S\ne V$. The average degree is $2|E|/|V|$, and the girth is the length of a shortest cycle, taken to be $+\infty$ for forests.
\paragraph{Statement recorded in the supplied source.}
For every positive real $d$, does there exist an integer $g$ such that every graph $G$ with at least two vertices, average degree less than $d$, and girth at least $g$ has a matching cut? \sref{3120}{2}
\subsection{Short English statement}
Does sufficiently large girth force a matching cut when the average degree has a fixed upper bound?
\subsection{Sources}
\begin{description}
\item[\hypertarget{source:3120:1}{S1}] ULAM.AI and the UnsolvedMath Contributors, UnsolvedMath: A Curated Collection of Open Mathematics Problems, record 3120 (OPG-37364). CC BY 4.0. \url{https://huggingface.co/datasets/ulamai/UnsolvedMath}.
\item[\hypertarget{source:3120:2}{S2}] Open Problem Garden, Matching cut and girth, original node 37364, statement and mathematical discussion. \url{https://www.openproblemgarden.org/op/matching_cut_and_girth}.
\end{description}
\label{end:3120}
\end{document}
